\documentclass[11pt]{article}
\usepackage[margin=1in]{geometry}
\usepackage{amsmath,amssymb,amsthm,mathtools}
\usepackage{graphicx}
\usepackage{xcolor}
\usepackage{microtype}
\usepackage[hidelinks]{hyperref}
\usepackage{enumitem}

\newtheorem{theorem}{Theorem}
\newtheorem{proposition}{Proposition}
\newtheorem{remark}{Remark}
\newcommand{\R}{\mathbb R}
\newcommand{\Z}{\mathbb Z}
\newcommand{\Sph}{\mathbb S^{d-1}}
\newcommand{\card}{\operatorname{card}}
\newcommand{\dist}{\operatorname{dist}}

\title{The translated-cone lattice count in every dimension}
\author{Solution packet for arXiv:2403.10350}
\date{13 August 2026}

\begin{document}
\maketitle

\begin{abstract}
Aksentijevi\'c, Aleksi\'c, and Pilipovi\'c ask whether two cones in good
position satisfy the translated-intersection bound needed for their product
theorem in dimension $d\ge3$, conjecturing exponent $d$.  We prove the bound,
with an explicit constant, whenever the closed spherical traces of the cones
are antipodally separated.  This includes the natural closed-cone
interpretation of their hypothesis and works in every dimension.  The exponent
$d$ is optimal.  We also resolve an important terminology ambiguity: if
``cone'' means open cone and only literal disjointness is imposed, the claim is
false even in dimension two, with one translated intersection containing
infinitely many lattice points.
\end{abstract}

\section{Source question}

For cones $\Gamma_1,\Gamma_2\subset\R^d\setminus\{0\}$, condition (a) of
Theorem 5.5 in the source asks for constants $C>0$ and $\gamma\ge1$ such that
\begin{equation}\label{eq:count}
 N(n):=\card\{k\in\Z^d:n-k\in\Gamma_1,\ k\in\Gamma_2\}
 \le C|n|^\gamma,\qquad n\in\Z^d.
\end{equation}
Remark 2 on page 13 of the arXiv PDF says that
$\Gamma_1\cap(-\Gamma_2)=\varnothing$ yields \eqref{eq:count} for $d=2$
with $\gamma=2$, and hypothesizes that it yields the bound for $d\ge3$ with
$\gamma=d$.  The page is reproduced in Figure~\ref{fig:source}.

The source uses ``open convex cone'' earlier, but its planar illustration in
Remark 2 uses non-strict inequalities.  The distinction matters.  We therefore
state the exact geometric hypothesis, prove the intended result, and then
treat the literal open-cone reading.

\section{Proof intuition}

Every counted $k$ gives a decomposition $n=x+y$ with
$x=n-k\in\Gamma_1$ and $y=k\in\Gamma_2$.  If the directions of $x$ and
$-y$ stay a positive angular distance apart, $x$ and $y$ cannot nearly cancel:
$|x+y|$ controls $|x|+|y|$.  Hence $|k|=|y|=O(|n|)$, so all admissible lattice
points lie in an $O(|n|)$ ball.  Counting the containing integer cube gives
$O(|n|^d)$ without describing any faces of the translated cone intersection.

\section{The all-dimensional theorem}

For a conic set $\Gamma\subset\R^d\setminus\{0\}$, write
\[
 K(\Gamma)=\overline{\Gamma\cap\Sph}^{\,\Sph}.
\]

\begin{theorem}[Sharp translated-cone count]\label{thm:main}
Let $d\ge1$ and let $\Gamma_1,\Gamma_2\subset\R^d\setminus\{0\}$ be
conic sets.  Assume
\begin{equation}\label{eq:sep}
 K(\Gamma_1)\cap K(-\Gamma_2)=\varnothing,
\end{equation}
and put
\[
 \delta=\dist\bigl(K(\Gamma_1),K(-\Gamma_2)\bigr)>0.
\]
Then $N(0)=0$ and, for every nonzero $n\in\Z^d$,
\begin{equation}\label{eq:explicit}
 N(n)\le \left(1+\frac{4|n|}{\delta}\right)^d
 \le \left(1+\frac4\delta\right)^d|n|^d.
\end{equation}
In particular, condition \eqref{eq:count} holds with $\gamma=d$.
If the two cones are closed relative to $\R^d\setminus\{0\}$, the source's
hypothesis $\Gamma_1\cap(-\Gamma_2)=\varnothing$ implies
\eqref{eq:sep}, so its high-dimensional conjecture follows.
\end{theorem}

\begin{proof}
If either cone is empty there is nothing to prove.  Otherwise the two sets in
\eqref{eq:sep} are disjoint compact subsets of the sphere, whence $\delta>0$.
Take nonzero $x\in\Gamma_1$ and $y\in\Gamma_2$, and set
\[
 a=|x|,\quad b=|y|,\quad u=x/a,\quad v=-y/b.
\]
Then $|u-v|\ge\delta$, and the elementary identity
\begin{align*}
 |x+y|^2
 &=|au-bv|^2\\
 &=(a-b)^2+ab|u-v|^2\\
 &\ge (a-b)^2+ab\delta^2
 \ge \frac{\delta^2}{4}(a+b)^2
\end{align*}
holds because $0<\delta\le2$ and
$(a-b)^2+4ab=(a+b)^2$.  Thus
\begin{equation}\label{eq:nocancel}
 |x+y|\ge\frac\delta2(|x|+|y|).
\end{equation}

For a point counted by $N(n)$, use $x=n-k$ and $y=k$ in
\eqref{eq:nocancel}.  It follows that
\[
 |k|\le |x|+|y|\le\frac{2|n|}{\delta}=:R.
\]
The lattice points in the radius-$R$ ball are contained in the cube
$[-R,R]^d$, so their number is at most
$(2\lfloor R\rfloor+1)^d\le(2R+1)^d$.  This is the first inequality in
\eqref{eq:explicit}; the second uses $|n|\ge1$.  When $n=0$, an admissible
$k$ would give $-k\in\Gamma_1\cap(-\Gamma_2)$, contradicting
\eqref{eq:sep}.

Finally, if the cones are relatively closed, their spherical traces are
already compact.  Their literal antipodal disjointness is therefore exactly
\eqref{eq:sep}.
\end{proof}

\begin{proposition}[Exact robust geometric condition]\label{prop:robust}
For arbitrary nonempty conic sets $\Gamma_1,\Gamma_2$, condition
\eqref{eq:sep} is equivalent to the existence of $c>0$ such that
\begin{equation}\label{eq:robust}
 |x+y|\ge c(|x|+|y|)
 \quad(x\in\Gamma_1,\ y\in\Gamma_2).
\end{equation}
One may take $c=\delta/2$.
\end{proposition}

\begin{proof}
The forward implication is \eqref{eq:nocancel}.  Conversely, if the closed
spherical traces meet at $u$, choose unit $u_j\in\Gamma_1$ and unit
$v_j\in-\Gamma_2$ tending to $u$.  Applying \eqref{eq:robust} to
$x=u_j$ and $y=-v_j$ gives $|u_j-v_j|\ge2c$, a contradiction.
\end{proof}

Thus closed-spherical antipodal separation is the full characterization of
the uniform no-cancellation mechanism behind the conjectured estimate.  It is
not asserted to classify exceptional arithmetic behavior of tangent cones.

\section{Why literal open-cone disjointness is insufficient}

\begin{theorem}[Open-cone counterexample]\label{thm:counter}
There are acute open convex cones $\Gamma_1,\Gamma_2\subset\R^2$ satisfying
$\Gamma_1\cap(-\Gamma_2)=\varnothing$ for which $N(n)=\infty$ at one fixed
$n\in\Z^2$.  There are full-dimensional open convex pointed examples in
every $d\ge3$ as well.
\end{theorem}

\begin{proof}
In dimension two set
\[
 \Gamma_1=\{(x_1,x_2):x_1>0,\ 0<x_2<x_1\},\qquad
 \Gamma_2=\{(x_1,x_2):x_1<0,\ 0<x_2<-x_1\}.
\]
Both are open convex cones of angular width $\pi/4$.  Every point of
$\Gamma_1$ has positive second coordinate, while every point of
$-\Gamma_2$ has negative second coordinate, so the required intersection is
empty.  Nevertheless, for $n=(0,2)$ and every integer $m\ge2$,
\[
 k_m=(-m,1)\in\Gamma_2,\qquad n-k_m=(m,1)\in\Gamma_1.
\]
Hence all the distinct $k_m$ are counted and $N(n)=\infty$.

For $d\ge3$, append the strict linear inequalities
\[
 |x_j|<\varepsilon x_2\quad(3\le j\le d)
\]
to each displayed cone, with a fixed sufficiently small $\varepsilon>0$.
The resulting sets are full-dimensional open convex pointed cones; the
vectors $(\pm m,1,0,\ldots,0)$ give the same infinite family.
\end{proof}

This also shows why passing from disjoint open cones to disjoint closed
spherical traces is necessary: the two cones above approach a common
antipodal boundary ray, allowing arbitrarily large cancellation.

\section{Optimality and consequence for the source}

The exponent $d$ is optimal.  Take $\Gamma_1=\Gamma_2$ equal to the open
positive orthant and $n=(N,\ldots,N)$, $N\ge2$.  Each
$k\in\{1,\ldots,N-1\}^d$ is counted, so
\[
 N(n)=(N-1)^d\asymp |n|^d.
\]

For relatively closed cones, Theorem~\ref{thm:main} supplies condition (a) of
the source's Theorem 5.5 with the conjectured and best possible value
$\gamma=d$.  The same conclusion holds more generally under
\eqref{eq:sep}.  This closes the geometric gap identified in Remark 2 and
makes the source's ensuing product corollary available in all dimensions,
subject to its remaining analytic hypotheses.  The packet does not
independently audit those other hypotheses.

\section{Verification and novelty limits}

The proof uses only compactness of closed subsets of the unit sphere, the
displayed exact norm identity, and the elementary integer-cube count.  The
counterexample is verified symbolically by its strict inequalities for every
$m\ge2$; no numerical approximation is involved.  Eight progressive proof
and stress-test attempts are recorded in the accompanying attempt log.

The repository's result, solution, attempt, and proof-gap indexes contained no
treatment of arXiv:2403.10350.  Bounded searches on 13 August 2026 used the
exact title, arXiv identifier, exact open-problem phrases, authors, and core
terms about translated cone lattice counts.  No later resolution was found.
The published 2024 Filomat paper retains the open statement, and a 2025
project review repeats it.  Since the decisive lemma is elementary geometry,
novelty should nevertheless be treated cautiously pending specialist citation
review.

Human review should first confirm the source's intended convention for
``cone.''  If cones are relatively closed, Theorem~\ref{thm:main} is a direct
full affirmative answer.  If they are open, Theorem~\ref{thm:counter} is a
counterexample to the literal claim and Theorem~\ref{thm:main} gives the sharp
statement repair.

\begin{figure}[p]
\centering
\includegraphics[width=0.96\textwidth]{figures/open_problem_crop.png}
\caption{Remark 2 on page 13 of the arXiv PDF: the dimension-$d$ cone-count
hypothesis, followed by the planar discussion.}
\label{fig:source}
\end{figure}

\begin{thebibliography}{9}
\bibitem{source}
A.~Aksentijevi\'c, S.~Aleksi\'c, and S.~Pilipovi\'c,
\emph{On the product of periodic distributions. Product in shift-invariant
spaces}, arXiv:2403.10350 (2024); Filomat \textbf{38} (2024), no.~23,
8011--8021.
\end{thebibliography}

\end{document}
